\chapter{Gaussian Vectors, Graphical Models, and Message Passing}
\label{ch:graphs}\label{ch:bp}

The research chapters compute diffusion scores by Bayesian inference on a
chain. This chapter builds the inference machinery itself, in the depth
its central role deserves, and in one place: the multivariate Gaussian
and its two parametrisations, the matrix tools (spectral theorem,
Neumann series) that the Gaussian analysis runs on, Markov random fields
and the Hammersley--Clifford theorem, factor graphs, the sum--product
algorithm and its exactness on trees, the Gaussian message family and its
classical incarnation, the Kalman filter, and the mean-field closures
(TAP, AMP) of statistical physics. The chapter closes with the
architectural side of the same story: convolutional networks, receptive
fields, and the recent identification of U-Nets with belief propagation
on hierarchical models. All of this material is standard; references
are given section by section, and the research chapters use every part
of it.

\section{The multivariate Gaussian, in both parametrisations}
\label{sec:gauss-vectors}

\begin{definition}[Multivariate Gaussian]\label{def:gaussian}
A random vector $a \in \R^K$ is Gaussian with mean $m$ and (positive
definite) covariance $\Sigma$, written $a \sim \Nd(m, \Sigma)$, if its
density is
\begin{equation}
  \Nd(a;\, m, \Sigma)
  \;=\; \frac{1}{(2\pi)^{K/2} \sqrt{\det \Sigma}}
  \exp\!\Big[ -\tfrac{1}{2} (a - m)^\T \Sigma^{-1} (a - m) \Big] .
  \label{eq:gauss-density}
\end{equation}
\end{definition}

Expanding the quadratic form shows that the same density can be written
in \emph{natural} (information) parameters,
\begin{equation}
  \Nd(a;\, m, \Sigma) \;\propto\;
  \exp\!\Big[ -\tfrac{1}{2}\, a^\T J\, a + h^\T a \Big],
  \qquad
  J = \Sigma^{-1},
  \qquad
  h = \Sigma^{-1} m ,
  \label{eq:gauss-natural}
\end{equation}
with $J$ the \emph{precision matrix} and $h$ the precision-weighted mean.
The two parametrisations are dual, and the duality organises everything
that follows \citep{bishop2006pattern, koller2009probabilistic}:

\begin{itemize}
  \item \emph{Linear maps.} If $a \sim \Nd(m, \Sigma)$ then
  $Aa + b \sim \Nd(Am + b,\, A \Sigma A^\T)$: Gaussianity is preserved by
  affine transformations, with moments transforming linearly. In
  particular a sum of independent Gaussians is Gaussian with means and
  covariances adding.
  \item \emph{Marginalisation is easy in $\Sigma$.} The marginal of a
  sub-block $a_1$ of $a = (a_1, a_2)$ is Gaussian with mean $m_1$ and
  covariance $\Sigma_{11}$: simply read off the block.
  \item \emph{Conditioning is easy in $J$.} Partitioning
  $J = \begin{psmallmatrix} J_{11} & J_{12} \\ J_{12}^\T & J_{22}
  \end{psmallmatrix}$, the conditional law of $a_1$ given $a_2$ is
  Gaussian with precision $J_{11}$ (read off the block) and mean shifted
  by $-J_{11}^{-1} J_{12}\, a_2$. Eliminating a variable instead produces
  the \emph{Schur complement} correction on its neighbours. Proofs of
  both statements are in Appendix~\ref{app:gauss}.
  \item \emph{Products are easy in $J$.} Multiplying two Gaussian
  densities in the same variable adds their natural parameters,
  $(J, h) \mapsto (J_1 + J_2,\, h_1 + h_2)$. Bayesian updating with
  Gaussian likelihoods is addition.
\end{itemize}

The dual reading of zeros is used constantly in this thesis. A
zero in the \emph{covariance} means marginal independence of two
coordinates. A zero in the \emph{precision} means \emph{conditional}
independence given all the others: from the conditioning rule, if
$J_{ij} = 0$ then $a_i$ and $a_j$ do not interact once the rest is fixed.
A Gaussian vector can be, and in this thesis always is, globally
correlated (dense $\Sigma$) while conditionally local (sparse $J$). Which
matrix one should look at depends on the question, and confusing the
two is a common source of error in reasoning about Gaussian dependence.

\section{Matrix tools: the spectral theorem and matrix inversion}
\label{sec:matrix-tools}

Since Gaussian computation is matrix computation, we record the two
linear-algebra facts the research chapters use constantly.

\begin{theorem}[Spectral theorem for symmetric matrices]
\label{thm:spectral}
Every symmetric $A \in \R^{K \times K}$ has an orthonormal basis of
eigenvectors: $A = U \operatorname{diag}(\omega_1, \dots, \omega_K)
U^\T$ with $U^\T U = I$ and real eigenvalues $\omega_i$.
\end{theorem}

Three consequences matter here. First, analytic operations act on
eigenvalues alone: if no eigenvalue vanishes,
\begin{equation*}
  A^{-1} = U \operatorname{diag}(1/\omega_i)\, U^\T,
  \qquad\text{and more generally}\qquad
  f(A) = U \operatorname{diag}\big(f(\omega_i)\big)\, U^\T .
\end{equation*} Second, affine images
share the eigenbasis: $(cA + dI)\, u_i = (c\,\omega_i + d)\, u_i$, so a
whole family of matrices of the form $c(t) A + d(t) I$ is diagonalised
once and for all by the $U$ of $A$. Chapter~\ref{ch:gaussian} uses
exactly this fact to diagonalise the diffusion problem at every noise
level simultaneously. Third, for the symmetric
Toeplitz covariances produced by stationary processes
(Section~\ref{sec:stationarity}), the eigenvectors are asymptotically the
Fourier modes of the sequence \citep{brockwell1991time}: stationary
Gaussian problems decouple, in the large-$K$ limit, into independent
scalar problems per frequency.

Inversion, when the inverse is not wanted eigenvalue by eigenvalue, has
two workhorse expansions. The \emph{Neumann series}
\begin{equation}
  (I - B)^{-1} \;=\; \sum_{m \geq 0} B^m ,
  \qquad \text{convergent for } \norm{B} < 1 ,
  \label{eq:neumann}
\end{equation}
expresses the inverse of a near-identity matrix as a sum of powers; it is
the tool behind the small-$t$ expansions of Chapter~\ref{ch:gaussian},
where powers of a banded matrix have growing but controlled bandwidth.
The \emph{Woodbury identity} relates the inverse of a matrix to the
inverse of a low-rank or structured perturbation of it; its statement,
proof, and the specific consequence used by the research chapters are
collected in Appendix~\ref{app:gauss}. Finally, for the tridiagonal
Toeplitz matrices that Markov chains produce, the inverse is available in
closed form through a transfer-matrix (geometric-ansatz) computation,
carried out where it is needed
(Chapter~\ref{ch:gaussian}).

\section{Markov random fields}\label{sec:mrf}

The Gaussian story above assigned structure to matrices. Graphical
models assign it to distributions, for arbitrary (not necessarily
Gaussian) variables.

\begin{definition}[Markov random field]
Let $G = (V, \mathcal{E})$ be an undirected graph with one random
variable $x_i$ per node. The collection $x = (x_i)_{i \in V}$ is a
\emph{Markov random field} (MRF) if it satisfies the local Markov
property: for every node $i$,
\begin{equation}
  p\big(x_i \,\big|\, x_{V \setminus \{i\}}\big)
  \;=\; p\big(x_i \,\big|\, x_{\partial i}\big),
  \label{eq:local-markov}
\end{equation}
where $\partial i$ is the set of neighbours of $i$ in $G$.
\end{definition}

The fundamental structure theorem connects this
conditional-independence definition to the energy picture of
Chapter~\ref{ch:ebm}: by the Hammersley--Clifford theorem
\citep[proved in][]{besag1974spatial}, a strictly positive density is an
MRF on $G$ if and only if it factorises over the \emph{cliques} of $G$,
\begin{equation}
  p(x) \;=\; \frac{1}{Z} \prod_{C \in \mathrm{cliques}(G)} \psi_C(x_C)
  \;=\; \frac{1}{Z}\, \exp\Big[ -\sum_C E_C(x_C) \Big] .
  \label{eq:hammersley-clifford}
\end{equation}
Conditional independence, graphical locality, and additive local
energies are one and the same thing. The Ising model \eqref{eq:ising} is
the textbook example. The decisive example for us is one-dimensional:
the chain factorisation \eqref{eq:chain-factorisation} makes a Markov
chain an MRF on a line, with pairwise nearest-neighbour energies. For
\emph{Gaussian} MRFs the correspondence becomes matrix-concrete: the
energy is the quadratic form $\tfrac12 a^\T J a - h^\T a$, its pairwise
terms are the entries $J_{ij}$, and the graph of the MRF is exactly the
sparsity pattern of the precision matrix. The dual reading of zeros of
Section~\ref{sec:gauss-vectors} and the Hammersley--Clifford theorem are
the same statement, once in linear algebra and once in probability. A
Gaussian Markov chain therefore has a tridiagonal precision, a fact
Chapter~\ref{ch:gaussian} derives explicitly and exploits throughout.

One more structural remark that the research chapters rely on: if each
variable of an MRF is observed through its own independent noisy channel,
the likelihood factors touch one variable each, so no clique is
enlarged. \emph{The posterior of an MRF under per-node observations is an
MRF on the same graph.} Observation changes the potentials, never the
graph.

\section{Factor graphs}\label{sec:factor-graphs}

MRFs record \emph{which} variables interact but not \emph{how} the joint
factorises: a triangle of pairwise terms and a single triple term have
the same graph. The representation that makes the factorisation itself
graphical, and on which message passing is most naturally defined, is
the \emph{factor graph} \citep{kschischang2001factor,
koller2009probabilistic}.

\begin{definition}[Factor graph]
Given a factorisation
$p(x) = \frac{1}{Z}\prod_{a} \psi_a(x_{\partial a})$, its factor graph is
the bipartite graph with a \emph{variable node} for each $x_i$, a
\emph{factor node} for each $\psi_a$, and an edge $(i, a)$ whenever
$x_i \in x_{\partial a}$.
\end{definition}

A connected factor graph with one fewer edge than nodes is a
\emph{tree}: it contains no cycles, so any edge separates it into two
disjoint components. Trees are the tractable case of inference: on
them, marginals and partition functions are computable exactly by the
recursion of the next section, whatever the alphabet of the variables.
The factor graph of a Markov chain observed through per-node noise is a
tree of a particularly simple shape (a spine of variables and transition
factors, with one evidence leaf per variable), as
Chapter~\ref{ch:gaussian} verifies by counting.

\section{The sum--product algorithm}\label{sec:sum-product}

On a factor graph, sum--product belief propagation (BP) passes messages
along every edge \citep{pearl1988probabilistic, kschischang2001factor,
mezard2009information}. Variable-to-factor messages multiply the
incoming factor messages from all \emph{other} neighbours, and
factor-to-variable messages integrate the factor against the incoming
variable messages:
\begin{equation}
  \nu_{i \to a}(x_i) \propto
    \prod_{b \in \partial i \setminus a} \hat\nu_{b \to i}(x_i),
  \qquad
  \hat\nu_{a \to i}(x_i) \propto
    \int \psi_a(x_{\partial a})
    \prod_{j \in \partial a \setminus i} \nu_{j \to a}(x_j)\;
    \dd x_{\partial a \setminus i},
  \label{eq:bp-rules}
\end{equation}
with the belief at node $i$ proportional to
$\prod_{a \in \partial i} \hat\nu_{a \to i}(x_i)$. The ``exclude the
recipient'' rule is the heart of the algorithm: the message $i \to a$
summarises everything node $i$ knows \emph{except} what it learned from
$a$, so that no piece of evidence is ever counted twice along a path.

\begin{theorem}[Tree exactness of sum--product]\label{thm:tree-exact}
On a tree-shaped factor graph, the equations \eqref{eq:bp-rules} have a
unique solution, reached after one inward and one outward pass, and the
resulting beliefs are the exact marginals of $p$
\citep[Theorem~14.1]{mezard2009information}.
\end{theorem}

The reason is the separation property of trees: each edge splits the
graph into two disjoint halves, and the message crossing it summarises
exactly the half behind it. On graphs with cycles the same updates,
iterated, define \emph{loopy} BP, which is extraordinarily useful but
approximate \citep{weiss2001correctness, malioutov2006walk}; none of that
is needed in this thesis, because our graphs are trees.

\paragraph{The representation problem.} For \emph{discrete} variables
each message is a finite probability vector and each update in
\eqref{eq:bp-rules} is finite arithmetic. For \emph{continuous} variables
each message is a probability \emph{density}, an infinite-dimensional
object, and the update integrals must be carried out in function space.
Tree exactness (a property of the graph) says nothing about
representability of the messages (a property of the model's function
class); keeping the two apart is essential, and the distinction is where
all the difficulty of the non-Gaussian research chapter lives. There are
three ways to make continuous BP computable: restrict to models whose
messages close in a finite-dimensional family (the Gaussian route, next
section); represent messages numerically on a grid (the route our
non-Gaussian reference solver takes, Chapter~\ref{ch:laplace}); or
project each message onto a tractable family and accept the error (the
route whose cost this thesis measures).

\section{Gaussian belief propagation and the Kalman filter}
\label{sec:gaussian-bp}

The Gaussian family is the canonical closed message family, and its
closure is pure algebra.

\begin{toolbox}[tb:closure]{The two closure lemmas of Gaussian message passing}
\emph{Lemma 1 (products).} For Gaussians in the same variable,
\begin{equation*}
  \Nd(a; m_1, v_1) \times \Nd(a; m_2, v_2)
  \;\propto\; \Nd(a; m, v),
  \qquad
  \frac{1}{v} = \frac{1}{v_1} + \frac{1}{v_2},
  \qquad
  \frac{m}{v} = \frac{m_1}{v_1} + \frac{m_2}{v_2} :
\end{equation*}
the sum of two quadratics in $a$ is a quadratic in $a$; precisions and
precision-weighted means add, as the natural-parameter rule of
Section~\ref{sec:gauss-vectors} requires.

\emph{Lemma 2 (affine marginalisation).} Pushing a Gaussian through a
linear-Gaussian transition,
\begin{equation*}
  \int \Nd(a'; \alpha a, \sigma^2)\; \Nd(a; m, v)\; \dd a
  \;=\; \Nd(a';\, \alpha m,\, \alpha^2 v + \sigma^2),
\end{equation*}
because $a' = \alpha a + \eta$ is a sum of independent Gaussians. Read in
the reverse direction the same computation gives
$\Nd\big(a;\, m'/\alpha,\, (v' + \sigma^2)/\alpha^2\big)$.

A model whose factors are Gaussian in each variable performs, in
\eqref{eq:bp-rules}, only these two operations. For such models every
message is Gaussian, carried exactly by two numbers, and each update is
a few arithmetic operations. \hfill$\square$
\end{toolbox}

The oldest and most famous instance predates factor graphs by decades.
For linear-Gaussian state-space models (a Gaussian Markov chain observed
through linear Gaussian noise), the forward sum--product sweep is the
\emph{Kalman filter} \citep{kalman1960new}, alternating a
\emph{predict} step (Lemma~2, push the belief through the dynamics) and
an \emph{update} step (Lemma~1, absorb the new observation, with the
correction weighted by the Kalman gain); the backward sweep combined with
the forward one yields the fixed-interval smoother of
\citet{rauch1965maximum}, which conditions each state on \emph{all}
observations. That the classical filter and smoother are exactly the
Gaussian specialisation of sum--product on the chain's factor graph was
made explicit by \citet{kschischang2001factor}; textbook treatments are
\citet{bishop2006pattern, koller2009probabilistic}.
Chapter~\ref{ch:gaussian} derives the filter from the BP sweep, line by
line, for our diffusion posterior, where the observation factors take a
specific form (a pseudo-observation whose strength depends on the
diffusion time), and verifies the equivalence numerically to
double-precision rounding.

One remark, recorded here because it resolves a doubt raised repeatedly
in supervision: the Gaussian closure of this section is \emph{algebra},
valid at any node degree, and invokes no limit theorem. It must not be
confused with the central-limit justification of the \emph{approximate}
Gaussian closures of the next section, which does require many incoming
messages and which a chain, with two neighbours per node, does not
provide.

\section{From BP to TAP and AMP}\label{sec:amp-intro}

Statistical physics developed its own message-passing calculus for
\emph{dense} graphs, where every node interacts weakly with many others.
On such graphs the incoming messages to a node are numerous and nearly
independent, so their product is governed by a central-limit argument:
the cavity field is approximately Gaussian, and each message can be
summarised by a mean and a variance even when the underlying variables
are not Gaussian. Historically the resulting self-consistency equations
appeared first as the Thouless--Anderson--Palmer (TAP) equations of the
Sherrington--Kirkpatrick spin glass \citep{thouless1977solution}, whose
distinctive ingredient is the Onsager reaction term, the correction that
removes a node's own influence from the field it feels. Three decades
later the same structure re-emerged in compressed sensing as
\emph{approximate message passing} (AMP) \citep{donoho2009message}: a
reduction of BP on a dense bipartite graph to iterations on means and
variances, with the Onsager term reappearing as a memory correction. The
derivation of AMP from BP by moment expansion is carried out for the
$\ell_2$ minimisation problem by \citet{genovese2026amp}; the general
technology, its state-evolution analysis, and its phase diagrams are
reviewed by \citet{zdeborova2016statistical} and, in lecture-note form,
by \citet{krzakala2024statphys}. \citet{mezard2017meanfield} derived the
TAP equations of the Hopfield model, connecting this calculus back to the
networks of Chapter~\ref{ch:statmech}.

The logic of the CLT justification should be kept explicit, because this
thesis probes exactly its boundary. On a dense graph with $N$ weak
interactions per node, each message removes one interaction out of $N$,
message and marginal differ at order $1/N$, and sums of many weak,
nearly independent terms are Gaussian: the per-node Gaussian closure is
controlled. On a chain every node has two neighbours: no aggregation
takes place, each message carries the entire influence of half the
system, and nothing justifies a Gaussian form for it beyond the special
cases where it is exact. What happens when one imposes the closure
anyway is a quantitative question. This thesis answers it on the chain
in two complementary ways: exactly, for the per-node (AMP/TAP) closure
on the Gaussian chain, where the breakdown admits a closed-form phase
boundary (Appendix~\ref{app:amp}); and numerically, for the
Gaussian-projected messages on genuinely non-Gaussian chains
(Chapter~\ref{ch:laplace}).

\section{Locality in learned architectures: convolutions, receptive
fields, U-Nets}\label{sec:cnn-locality}

The message-passing story has an architectural mirror, and it supplies
the vocabulary for the locality questions of the research chapters. A
convolutional network processes a signal by applying the same local
filter at every position, stacking such layers so that depth enlarges
the \emph{receptive field}: the window of input coordinates that can
influence a given output coordinate \citep{lecun2015deep}. Convolutional
architectures powered the breakthrough of deep learning in vision
\citep{krizhevsky2012imagenet}, and their central inductive bias is
precisely locality plus translation invariance: the value of an output
pixel is computed from a bounded neighbourhood, identically across the
image. The standard denoising backbone of diffusion models, the U-Net
\citep{ronneberger2015unet, ho2020denoising}, is a convolutional
encoder--decoder with skip connections, whose receptive field grows
along the contracting path and whose skip connections reinject
fine-scale local information along the expanding path.

Why local architectures should be \emph{sufficient} for score estimation
is a theoretical question, and message passing has begun to answer it.
\citet{mei2024unets} shows that for generative hierarchical models
(tree-structured graphical models over image-like data), the denoising
function computed by belief propagation can be implemented efficiently
by a U-Net, with the network's layers mirroring the sweeps of the
message-passing schedule; convolutional layers are efficient
implementations of the local BP updates. This gives the U-Net an
inference-algorithm reading: the architecture is good at denoising
because denoising, on structured data, \emph{is} message passing, and
message passing is local.

The present thesis works one step below this correspondence, on a model
small enough that the locality of the score is not an architectural
hypothesis but a computable quantity. For the Gaussian chain, the exact
score at frame $k$ depends on the observation at frame $k \pm d$ with a
weight that decays geometrically in $d$, at a rate $q(\alpha, t)$
computed in closed form in Chapter~\ref{ch:gaussian}; the error of the
best radius-$r$ local approximation decays as $q^r$, and the comparison
between truncating the inference and truncating the score matrix
quantifies, in one solvable case, how a local computation should be
organised. We state the obvious caveat once: no network is trained in
this thesis, so these results describe the function a network would have
to learn, not the network that learns it.
